\chapter{Taxonomy Tuning Protocol and Ablations}\label{app:tuning-protocol}

This appendix details the mechanics behind the canonical hyperparameters
and construction ablations referenced in
Chapter~\ref{ch:experimental-design}; Table~\ref{tab:hyperparams}
lists the selected values.

\begin{table}[htbp]
\centering
\caption{Canonical taxonomy construction hyperparameters.}
\label{tab:hyperparams}
\begin{tabular}{lll}
\toprule
Parameter & Symbol/Name & Value \\
\midrule
Embedding dimension & $d$ & 256 (fixed) \\
Maximum recursion depth & $D_{\max}$ & 8 \\
Minimum cluster size & $n_{\min}$ (\texttt{minClusterSize}) & 55 \\
Proposal separation bar & $\epsilon_{\text{sep}}$ (\texttt{proposalSeparationBar}) & 0.025 \\
Acceptance tolerance & $\tau$ (\texttt{tau}) & $10^{-6}$ \\
Acceptance gate, in paired SE units & \texttt{acceptanceZ} & 2.0 \\
Membership share floor\textsuperscript{a} & \texttt{membershipFloor} & 0.25 \\
Routing beam margin & \texttt{routingBeamGamma} & 0.20 \\
Descent-gate slack & \texttt{descentMargin} & 0.12 \\
Effective support floor & \texttt{effectiveSupportFloor} & 2.0 \\
Max leaf assignments (arena-time cap) & \texttt{maxLeafAssignments} & 5 \\
Fusion pre-filter threshold & \texttt{fusionSimilarityThreshold} & 0.90 \\
Iteration budget & \texttt{numIterations} & 35\textsuperscript{c} \\
Construction mode\textsuperscript{b} & \texttt{dagMode} & \texttt{DAG\_MAX} \\
\bottomrule
\end{tabular}
\par\smallskip
\footnotesize \textsuperscript{a}Minimum share of a query's own
normalised membership a reached destination must hold to count; the
descent decision and sibling admission are governed separately by the
descent gate and the beam margin
(Appendix~\ref{app:dag-formalism}).
\textsuperscript{b}The mode name is historical. It now selects three
flags (residual routing, the residual split gate, and stable question
identifiers) and carries no polyhierarchy setting, since cross-linking
is no longer part of the pipeline (Section~\ref{sec:poly-negative}).
\textsuperscript{c}An upper budget, not a run length: the frozen
construction reached its certified fixed point at \textbf{iteration
10} and iterations 8--10 are identical on every recorded column
(Section~\ref{sec:res-construction}).
Values are those of the frozen construction run that every result in
this thesis is attributed to, read from its run configuration.
Two of them are commonly stated backwards and are stated here in the
direction the evidence supports.
\texttt{minClusterSize} is a \textbf{budget} choice --- it sets the
cell count to what the judge budget can identify, and per-cell
reliability follows from it rather than motivating it --- and it is
free to choose only because granularity turned out to be evaluatively
flat (Section~\ref{sec:res-granularity}).
\texttt{descentMargin} is a \textbf{mode selector} with two justified
operating points rather than a tuned optimum: $0.12$ is
routing-optimal with sub-domain discovery off, which is what this
artifact reports, while a value near $0$ turns discovery on at a cost
of roughly five points of routing quality and a large residual
reserve.
\end{table}

\section{Hyperparameter Selection Protocol}\label{sec:exp-tuning-protocol}

The original L9($3^4$) orthogonal-array sweep documented here is
\textbf{historical}: it was conducted on a predecessor
parameterisation of the routing layer, and three of its four
factors (\texttt{routingSoftmaxTau}, \texttt{assignmentCosineGap},
and \texttt{deltaAssign}) were subsequently removed when routing was
redesigned around the descent gate, beam margin, and
membership-share floor (Appendix~\ref{app:dag-formalism}).
The current calibration sweep uses the same L9/Pareto machinery over
the four live factors
\texttt{proposalSeparationBar} $\times$ \texttt{descentMargin} $\times$
\texttt{minClusterSize} $\times$ \texttt{membershipFloor};
because construction is deterministic at a fixed seed
(Appendix~\ref{app:dag-formalism}), screening comparisons across
configurations are exact rather than confounded by run-to-run
variance.
Its protocol and gates are retained because they informed the
surviving canonical values (notably \texttt{minClusterSize}) and
because the selection machinery will be reused for the tuning study
of the current parameter set.
The historical design covered four factors, each at three levels, run
against a Pareto selection procedure. Only one of those factors
(\texttt{minClusterSize} $\in \{50, 75, 100\}$) survives in the current
parameterisation; the other three controlled a routing temperature and
two nat-margin assignment gates that were removed with the mechanism
they governed, and their levels are not reproduced here because no
current parameter corresponds to them.

Configuration selection proceeds in three steps. Every candidate is
first checked against the eight hard gates in
Table~\ref{tab:pareto-gates}; a single failure excludes it. The same
candidate is then scored against the six soft gates in
Table~\ref{tab:pareto-gates}, whose outcome is recorded per candidate
but does not exclude it. All hard-gate survivors are then ranked by
Pareto dominance count over twelve coherence and calibration
objectives: five maximised (Weighted Leaf Purity, Dendrogram Purity,
Spherical Silhouette, total migration volume $m_{\text{total}}$ (the
fraction of queries whose anchor assignment changed relative to the
canonical partition; \emph{not} the RQ2 fidelity metric $\Delta\rho$,
which no tuning gate ever observes),
canonical--adapted Jaccard) and seven minimised (Total Dasgupta Cost,
Routing ECE, Brier score, no-match rate, small-leaf fraction,
$\kappa$-shrinkage mean, selected-node starved-leaf
fraction). Raw Top-1 accuracy is the tiebreaker. Small-leaf
fraction, starved-leaf fraction, and canonical--adapted Jaccard
appear in both layers: the gate sets a hard floor, and the dominance
ranking additionally credits margin beyond it.

\begin{table}[htbp]
\centering
\caption{Configuration-selection gates used by the L9/Pareto tuning protocol.}
\label{tab:pareto-gates}
\begin{tabular}{p{5.4cm}p{3.6cm}}
\toprule
\multicolumn{2}{l}{\textit{Hard gates --- any failure excludes the candidate}} \\
\midrule
Acyclic & \texttt{true} \\
Multi-parent nodes & $0$ \\
Orphaned nodes & $0$ \\
Assignment-cap rate & $\leq 0.20$ \\
Small-leaf fraction & $\leq 0.50$ \\
Starved-leaf fraction (selected node) & $\leq 0.20$ \\
\midrule
\multicolumn{2}{l}{\textit{Soft gates --- recorded, not excluding}} \\
\midrule
Average match count & $[1.05,\,1.30]$ \\
Top-1 routing accuracy & $\geq 75.6\%$ \\
Routing ECE\textsuperscript{*} & $\leq 0.25$ \\
Borderline rate & $[0.20,\,0.35]$ \\
Cross-anchor migration rate & $[0.10,\,0.30]$ \\
Canonical--adapted Jaccard & $[0.40,\,0.70]$ \\
\bottomrule
\end{tabular}
\par\smallskip
\footnotesize \textsuperscript{*}The frozen run exported a routing ECE
of $0.2114$, inside this gate. The exported quantity is not the metric
Appendix~\ref{app:metrics} defines: the export aggregates leaf
membership onto domain ancestors by taking a maximum over leaves rather
than summing, and a maximum over leaves is not a distribution. It is
reported here as a routing diagnostic with the discrepancy named, and
nothing in this thesis depends on it
(Section~\ref{sec:disc-debt}).
\end{table}

\begin{table}[htbp]
\centering
\caption{Parameters locked across the L9 sweep and ablations.}
\label{tab:tab-locked}
\begin{tabular}{ll}
\toprule
Locked parameter & Value \\
\midrule
$D_{\max}$ & 8 \\
\texttt{effectiveSupportFloor} & 2.0 \\
\texttt{fusionSimilarityThreshold} & 0.90 \\
\texttt{defaultKappaPrior} & 10.0 \\
\texttt{maxLeafAssignments} & 5 \\
\texttt{dagMode} & \texttt{DAG\_MAX} \\
\texttt{numIterations} & 35 \\
$d$ & 256 (except A1) \\
\bottomrule
\end{tabular}
\end{table}

The dynamic-temperature exponent $\gamma$ was \textbf{not} part of the
L9 design: because $\tau_{\text{eff}} = \tau\kappa^{\gamma}$ coupled
non-linearly to $\kappa$, it was ablated as a standalone factor at
$\{0, 0.5, 1.0\}$ (A11) rather than orthogonally; like the other
temperature knobs, $\gamma$ has since been removed along with the
mechanism it controlled, so A11 is historical. No reserved-arena-set
label, accuracy, or judge verdict is used by either the hard or soft
gates at any point in the L9/Pareto procedure; this is the mechanism
by which the tuning/test-leakage boundary is enforced in practice, not
merely asserted.

Figure~\ref{fig:calibration-sweep} summarises the sensitivity of the
headline construction metrics to the four live calibration factors.

\begin{figure}[htbp]
\centering
\fbox{\parbox{0.9\textwidth}{\centering\vspace{0.6cm}
Figure A.1 --- Calibration sweep small multiples (placeholder)\\
\footnotesize Dot-plot grid: four factor columns (\texttt{proposalSeparationBar}, \texttt{descentMargin}, \texttt{minClusterSize}, \texttt{membershipFloor}) $\times$ metric rows (Weighted Leaf Purity, Routing ECE, leaf count); data: \texttt{tuning/combined\_ledger.csv}.\vspace{0.6cm}}}
\caption{Calibration sweep small multiples. Panels form a grid of the
four live factors (\texttt{proposalSeparationBar}, \texttt{descentMargin},
\texttt{minClusterSize}, \texttt{membershipFloor}) against
headline construction metrics (Weighted Leaf Purity, Routing ECE,
leaf count); each panel plots the metric against the factor's
levels with one dot per run and seed encoded by marker shade. Flat
panels indicate robustness to the factor; sloped panels identify the
sensitivities that determined the canonical values of
Table~\ref{tab:hyperparams}.}
\label{fig:calibration-sweep}
\end{figure}

\section{Ablation Matrix}\label{sec:appendix-ablations}

Table~\ref{tab:ablations} summarises the ablation experiments run
against the canonical configuration, each varying exactly one
hyperparameter or structural choice while holding all others fixed.

\begin{table}[htbp]
\centering
\caption{Taxonomy construction ablation matrix.}
\label{tab:ablations}
\begin{tabular}{p{1.0cm}p{4.2cm}p{7.2cm}}
\toprule
ID & Factor & Levels tested \\
\midrule
A1 & Embedding dimension $d$ & $\{128, 256, 512\}$; canonical $= 256$ \\
A2 & $\kappa$ estimator corrections (Newton--Raphson refinement, Hornik--Gr\"un shrinkage, parent-anchored EB blend) & On (canonical) vs. off \\
A3$^\dagger$ & NMI variant & Overlapping vs. disjoint Shannon entropy \\
A4$^\dagger$ & Purity formula & Multi-membership-aware vs. primary-assignment purity \\
A5 & Separation threshold $\epsilon_{\text{sep}}$ & $\{0.01, 0.025, 0.05, 0.10, 0.30\}$; canonical $= 0.025$ \\
A5b & Marginal separation requirement & $\{0, \text{fall back to } \epsilon_{\text{sep}}\}$; $0$ measured and rejected \\
A6$^\ddagger$ & \texttt{deltaAssign} $\times$ \texttt{assignmentCosineGap} (joint; historical) & $\{0.5,1.0,2.0\} \times \{0.10,0.15,0.20\}$ \\
A7 & Iteration count & $\{10, 25, 50\}$ \\
A8 & Query set scope & Reserved (arena-excluded) vs. full corpus \\
A9 & Soft-membership refitting \texttt{refitMuPerIteration} & On (canonical) vs. off \\
A10$^\ddagger$ & Cross-link insertion (historical) & Reported in Section~\ref{sec:poly-negative} \\
A11$^\ddagger$ & Dynamic-temperature exponent $\gamma$ (historical) & $\{0.0, 0.5, 1.0\}$; canonical $= 0.0$ \\
\bottomrule
\end{tabular}
\end{table}

\textbf{Status of these rows.} Construction is deterministic at a fixed
seed, so each ablation is a single run at seed 42 rather than a mean
over seeds; no ablation figure in this thesis is a multi-seed mean.
Three rows have current, citable outcomes and the rest do not.
The separation-threshold row (A5) is what fixed the canonical
$\epsilon_{\text{sep}} = 0.025$, positioned between the two nulls of
Section~\ref{sec:arch-requirements}.
The marginal-separation variant --- requiring cluster $k+1$ to add a
fixed amount over $k$ rather than falling back to the level bar --- was
measured at $0$ and \textbf{rejected}: 83 leaves at $J = 0.230247$
against the canonical artifact's 87 leaves at $J = 0.253129$, and it
relocates the binding threshold onto the maximum branching factor.
The maximum-branching-factor plateau and the $k$-fallback figures were
measured before a correction to the splitter's routing check and are
void as to values; the qualitative finding they support --- that the
branching cap was a structural determinant of the tree before the
fallback existed, and is a cost bound afterwards --- stands, and the
numbers are not restated.

A3 and A4 (marked $\dagger$) are not
system ablations in the same sense as the others: they vary the
\textit{metric computation formula} used to score a fixed,
already-constructed taxonomy, rather than a system parameter that
changes the taxonomy itself. They are retained as
\textit{reporting-sensitivity analyses}, testing whether headline
quality conclusions are an artefact of a specific metric definition.
A6, A10, and A11 (marked $\ddagger$) are \textit{historical}: they
ablated the two nat-margin secondary-assignment gates, the cross-link
insertion operator, and the dynamic routing temperature, all of which
have since been removed from the pipeline.
Their rows are retained for provenance; an equivalent ablation of
the current parameter set is pending, and the cross-link result is
reported in Section~\ref{sec:poly-negative}.
% TODO(calibration): finalize once taxonomy construction is frozen
A9 isolates soft-membership refitting, whose marginal contribution to
taxonomy quality needs to be demonstrated, not merely described.


\section{Pre-Registrations}\label{app:preregs}

Four decision rules were fixed in dated documents before the data they
govern existed. This section records what each fixed and when, so that
every use of the word ``pre-registered'' in
Chapter~\ref{ch:results} can be checked. The dates, from the
version-control history: P1, P2 and P3 were registered on 2026-07-27;
P4 was registered on 2026-07-27 and amended through 2026-07-28, with
every amendment predating the Run~C results it governs; P2 carries a
dated results-written-back addendum of 2026-07-28. The purpose of reproducing
them is not ceremony: three results in this project inverted after the
fact, and only the ones whose reading had been fixed in advance
survived the inversion as findings rather than as revisions.

\subsection*{P1 --- Granularity and discriminative power}

Fixed in advance: three named outcomes, so that any result the analysis
could return already had an interpretation attached; and the rule that
the \emph{observed} between-cell correlation is primary and the
disattenuated one secondary. The outcome that occurred was named in
advance as ``the same signal, thinner''. The primacy rule proved
decisive: the disattenuated column moves in the opposite direction
across granularities and clips past $1.0$, so leading with it would
have produced a conclusion from the arithmetic of the correction rather
than from the data (Section~\ref{sec:res-granularity}).

\subsection*{P2 --- Rubric specificity}

Fixed in advance: four directional predictions, each with the direction
of its own confound named separately --- including one measure whose
vocabulary-size confound could plausibly have driven the \emph{null}
arm higher, which is why it is reported alongside three others rather
than alone. Also fixed: within-arm decision thresholds, and the rule
that if the sharpest discriminator came back at zero in both arms the
test is \emph{inconclusive} rather than negative
(Section~\ref{sec:res-rubric-null}).

As originally registered the document recorded ``treatment arm only,
no null, this is not yet evidence''; a dated addendum (2026-07-28) has
since written the null-arm results back into it, resolved which
decision band governs which statistic (the within-arm band governs the
treatment arm's specificity; the $\leq 0.05$ / $\geq 0.10$ band of the
arena-launch document governs the random arm's absolute specificity),
and recorded that the null-arm artifacts are untracked in version
control, with the reproduction of the headline separation dated
2026-07-27 (Section~\ref{sec:res-rubric-null}).

\subsection*{P3 --- Arena launch cost}

Fixed in advance: cost bands for per-cell identification, with the
instruction that no band is to be renegotiated once the number is seen;
and a void condition --- if the measured identification cost came back
constant across cells, that would indicate the scheduler was still
capping and the measurement would report nothing. The measured cost
varies across cells and the minimum falls inside the pre-registered
band (Section~\ref{sec:res-cost}).

\subsection*{P4 --- The no-partition baseline}

Fixed in advance: the paired design (both arms judge an identical
question set, unlike the original MT-Bench protocol, because here the
partition is the treatment); the primary outcome (Spearman $\rho$
against ground-truth accuracy per domain) and the secondary
(per-verdict answer-key agreement, which has more power but measures
the judge rather than the grouping); the decisive difference of two
Spearman grid steps on Math, i.e.\ $|\Delta\rho| \geq 0.007$ at twelve
models; the designation of Math as the \emph{null} arm and law as the
test, with the reason; and four void conditions --- identical question
sets, the confidence gate disabled in both arms, identical query
samplers, and identically defined ties across the two verdict formats.

Recorded in the same document, before the run: the honest prior that
with the multiple-choice options present the comparison would likely
come back null, because the judge can verify the answer and the rubric
is then decorative. Writing down the expected outcome, with its
mechanism, in a dated document is what distinguishes the null of
Section~\ref{sec:res-c5} from an absence of effort.

\subsection*{P5 --- The options-blind arm, predicted but not run}

The options-blind arm has not been run
(Section~\ref{sec:disc-future}). Its predictions were nevertheless
fixed in writing, per stratum, and are reproduced because a prediction
written before a run that never happened is still evidence of
discipline --- and because it names its own alternative outcome.

Predicted: the untraced-against-untraced stratum, at $92.9\%$ with
options present, should return to chance with a very high tie rate;
if it comes back at $70\%$ or above, something is leaking, and the
candidates are enumerable --- the judge recognising the question from
pre-training, position bias surviving the dual-call control, or
answer-letter frequency priors in the corpus.
The traced-against-traced stratum, at $84.2\%$ with a standard error of
$2.5\%$, is the real measurement; a fall to about $75\%$ would be
$3.6$ standard errors and clearly readable.
\textbf{And the second possibility, written down in advance:} landing
near $84\%$ unchanged is plausible, and would mean that on traced pairs
the judge was never using the key at all, so the $88.8\%$ aggregate is
carried entirely by the untraced and mixed strata. That is a cleaner
decomposition than any currently available, and it should not be
discovered after the fact.

\section{Register of Corrections to Previously Reported Figures}\label{app:corrections-register}

Earlier versions of this pipeline and of its analysis record reported
figures that this thesis corrects or withdraws. Each correction is
told once, here; the main text states the corrected fact where the
number is used and points to this register for the story. No value
computed before a correction is quoted anywhere in the thesis.

\subsection*{Bradley--Terry standard errors}

The zero-sum rank completion requires subtracting $1/K^{2}$ after
inversion; the implementation subtracted $1/K$
(derivation: Appendix~\ref{app:numerics}). The over-subtraction of
$(K-1)/K^{2}$ --- $0.109$ at $K = 8$ --- exceeds the variance itself
at every realistic comparison count, so every Bradley--Terry standard
error this system reported before the correction was clamped at the
floor constant of $0.01$ rather than measured. Correct values are 44
to 209 times larger. This invalidates every confidence interval and
every ``resolved pair'' determination computed before the fix, since
pair resolution thresholds on $3(\sigma_i + \sigma_j)$.

\subsection*{The reliability constant, and the disattenuation form}

The constant in $r = n/(n+c)$ was published at $c = 7.66$. Two
independent errors compound. The curve was fitted against raw
split-half correlations, which measure reliability at \emph{half}
length, so fitting $n/(n + 2c_0)$ against $n/(n + c)$ returns
$c = 2c_0$; correcting the length mismatch halves it. And $c$ is a
property of the roster rather than of the corpus, and re-fitting at
the corrected 11-model band multiplies it by a further $0.66$. The
corrected value is $c = 2.52$. Separately, disattenuation must be
two-sided: both sides of $\rho(\text{arena}, \text{ground truth})$ are
finite-sample per-cell rankings, so the correction is $\rho / r$
rather than $\rho / \sqrt{r}$.

The consequences invert rather than shift. At a $40.8$-query cell an
observed median $\rho$ of $0.823$ corrects to $0.978$ under the
published constant and to $0.874$ under the corrected one. The
reliability ladder across granularities compresses from a spread of
$0.219$ to $0.088$, and leaf-scale cells stop being individually
low-reliability. That weakens the reliability leg of the argument for
judging at a coarser cut; the identification-cost leg of
Section~\ref{sec:res-cost} is untouched and probably carries it alone.
Two caveats travel with any value of $c$: the bootstrap interval
originally quoted for it resamples three fitted points and does not
propagate sampling uncertainty in the underlying split-half estimates,
so it is not a confidence interval; and the curve was fitted on
ground-truth rankings, which are noise-free per question, so arena
reliability at a given $n$ lies \emph{below} it. It is a ceiling, not
a prediction.

\subsection*{Evaluative-redundancy percentages}

Figures of the form ``$72.2\%$ of leaf pairs redundant, $42.2\%$
saturated'' were reported by an earlier version of this analysis. They
could not be reproduced under either an observed-primary or a
disattenuated reading, and the convention that produced them is
underspecified in the pre-registration. Under the corrected constant a
redundancy claim would in any case require an observed $\rho$ of
$0.844$ rather than $0.749$ to clear the same threshold, so the claim
would likely become a minority one. No redundancy percentage is
asserted in this thesis.

\subsection*{The correctness-blind rank correlation}

The value originally recorded for rank fidelity on the subset where
the answer key cannot discriminate --- Spearman $\rho$ falling from
the $0.95$--$0.98$ region to $0.74$--$0.76$ --- could not be
reproduced; recomputation on the same verdict set returns
$0.90$--$0.93$ under both tie conventions. The two computations differ
by an analysis convention that was not written down at the time. The
magnitude is withdrawn; the direction and the tie behaviour carry the
claim (Section~\ref{sec:res-correctness}).

\todo[inline]{Unresolved: the correctness-blind Spearman $\rho$.
Recover the convention that produced the $0.74$--$0.76$ figure (tie
weighting, row subsetting, and whether $\rho$ was fitted on
Bradley--Terry $\theta$ or on win-rate), then report one value per tie
convention or drop the magnitude claim entirely.}

\subsection*{Routing expected calibration error}

Routing ECE is implemented and was measured at $0.2114$ on the frozen
run, inside the gate set for it. But the export aggregates leaf
membership onto domain ancestors by taking a maximum over leaves
rather than a sum, and a maximum over leaves is not a distribution;
the exported number is therefore not the metric its definition
specifies. It is reported as a routing diagnostic with the discrepancy
named (Table~\ref{tab:pareto-gates}), and no claim rests on it.

\subsection*{A side-channel found before it could enter a run}

Recorded here because it belongs to the same audit trail, although no
reported figure was affected. A pre-launch format check on two
candidate roster entries had exactly one anomalous column: newline
count, with a median of $0$ where every other model showed 8 to 24.
The two systems' stored output was not prose but a raw JSON envelope
carrying a self-reported \texttt{reason\_code} field alongside the
reasoning, and that field predicts the very outcome being judged ---
$83.2\%$ correct at code A against $36.5\%$ at code C for one system,
$90.0\%$ against $57.5\%$ for the other. The rendering path returns
stored output verbatim, so a judge would have received the code
inline. An asymmetric side-channel present on two models and absent
from the rest is strictly worse than a format artifact; neither system
entered any arena run as stored.

\todo[inline]{Resolve before submission: (i) recompute every
Bradley--Terry interval quoted in Chapter~\ref{ch:results} against the
corrected variance term, from the stored verdicts; (ii) fix the
routing ECE aggregation (maximum $\rightarrow$ sum) and re-export, or
keep the current number with the discrepancy stated as it is here;
(iii) re-fit the reliability constant against the roster the reported
arena runs actually used, since $c = 2.52$ is the 11-model-band value
fitted on ground-truth rankings.}
